% Tệp thành phần: không chứa preamble hoặc môi trường document.
\section{Kiểm định giả thuyết hai mẫu độc lập}
\label{sec:two-sample-method}

Nhãn nhị phân của nghiên cứu gộp DoS, Probe, R2L và U2R thành nhóm Attack, trong
khi giữ Normal là nhóm đối chứng. Vì vậy, kiểm định hai mẫu được dùng để đánh giá
liệu trung bình của một đặc trưng lưu lượng có khác nhau giữa hai nhãn này hay không.
Phép kiểm định bổ sung cho ANOVA năm nhóm: ANOVA bảo toàn sự không đồng nhất giữa
bốn họ tấn công, còn kiểm định hai mẫu trả lời tương phản trực tiếp Normal--Attack.
Nó cũng khác Pearson, vì Pearson đo liên hệ tuyến tính giữa hai biến định lượng chứ
không trực tiếp kiểm định chênh lệch trung bình theo nhãn. Thống kê mô tả cho biết vị
trí và độ phân tán quan sát được, nhưng chưa định lượng độ bất định của chênh lệch.

Mục tiêu của phân tích là cung cấp bằng chứng suy diễn và kích thước khác biệt cho
những đặc trưng có ý nghĩa hành vi mạng. Kết quả không được xem là thước đo trực tiếp
của năng lực dự đoán và không chứng minh nhãn tấn công gây ra giá trị đặc trưng.

\subsection{Mục tiêu và câu hỏi nghiên cứu}
\label{subsec:two-sample-rq}

Ba câu hỏi được xác định trước khi xem kết quả kiểm định:
\begin{enumerate}[label=\textbf{RQ\arabic*:},leftmargin=1.7cm]
  \item Tỷ lệ lỗi SYN trung bình của các kết nối đến cùng dịch vụ
  (\path{srv_serror_rate}) có khác nhau giữa Normal và Attack hay không?
  \item Tỷ lệ lỗi REJ trung bình của các kết nối đến cùng dịch vụ
  (\path{srv_rerror_rate}) có khác nhau giữa hai nhóm hay không?
  \item Số kết nối gần đây trung bình (\texttt{count}) có khác nhau giữa hai nhóm
  hay không?
\end{enumerate}

Hai biến đầu là tỷ lệ trong miền $[0,1]$, lần lượt phản ánh lỗi SYN và lỗi REJ trong
tập kết nối đến cùng dịch vụ. \texttt{count} đếm số kết nối đến cùng máy đích trong
cửa sổ hai giây, nên là biến đếm rời rạc. Ba biến được chọn vì đại diện cho lỗi thiết
lập kết nối, lỗi từ chối và cường độ kết nối; chúng không phải một bước sàng lọc toàn
bộ 41 đặc trưng của NSL-KDD.

\subsection{Cơ sở lý thuyết của kiểm định hai mẫu độc lập}
\label{subsec:two-sample-theory}

Gọi $g\in\{N,A\}$ lần lượt là Normal và Attack; $n_g$, $\bar{x}_g$, $s_g^2$ và
$\mu_g$ là cỡ mẫu, trung bình mẫu, phương sai mẫu và trung bình tổng thể. Các thống
kê trong từng nhóm được xác định bởi
\begin{equation}
  \bar{x}_g=\frac{1}{n_g}\sum_{i=1}^{n_g}x_{gi},
  \label{eq:two-sample-mean}
\end{equation}
\begin{equation}
  s_g^2=\frac{1}{n_g-1}\sum_{i=1}^{n_g}(x_{gi}-\bar{x}_g)^2.
  \label{eq:two-sample-variance}
\end{equation}
Ước lượng điểm của chênh lệch trung bình là
\begin{equation}
  \widehat{\Delta}=\bar{x}_N-\bar{x}_A,
  \label{eq:two-sample-difference}
\end{equation}
trong đó dấu âm nghĩa là trung bình Attack lớn hơn trung bình Normal. Do phương sai
hai nhóm không được giả định bằng nhau, sai số chuẩn Welch là
\begin{equation}
  SE(\widehat{\Delta})=
  \sqrt{\frac{s_N^2}{n_N}+\frac{s_A^2}{n_A}}.
  \label{eq:welch-se}
\end{equation}
Thống kê kiểm định và bậc tự do Welch--Satterthwaite lần lượt là
\begin{equation}
  t=\frac{\bar{x}_N-\bar{x}_A}
  {\sqrt{\frac{s_N^2}{n_N}+\frac{s_A^2}{n_A}}},
  \label{eq:welch-t}
\end{equation}
\begin{equation}
  \nu=
  \frac{\left(\frac{s_N^2}{n_N}+\frac{s_A^2}{n_A}\right)^2}
  {\frac{\left(\frac{s_N^2}{n_N}\right)^2}{n_N-1}
  +\frac{\left(\frac{s_A^2}{n_A}\right)^2}{n_A-1}}.
  \label{eq:welch-df}
\end{equation}
Khoảng tin cậy hai phía cho chênh lệch trung bình được tính nhất quán với kiểm định:
\begin{equation}
  CI_{95\%}(\mu_N-\mu_A)=
  \widehat{\Delta}\pm t_{0.975,\nu}SE(\widehat{\Delta}).
  \label{eq:welch-ci}
\end{equation}
Các khoảng ở Công thức~\eqref{eq:welch-ci} là khoảng riêng lẻ cho từng RQ, không tự
động bảo đảm độ bao phủ đồng thời 95\% cho cả ba câu hỏi.

Ngoài chênh lệch trên đơn vị gốc, nghiên cứu báo cáo Cohen $d$ và Hedges $g$ theo quy
ước độ lệch chuẩn gộp. Với $s_p$ là độ lệch chuẩn gộp,
\begin{equation}
  d=\frac{\bar{x}_N-\bar{x}_A}{s_p},\qquad
  g=\left(1-\frac{3}{4(n_N+n_A)-9}\right)d.
  \label{eq:hedges-g}
\end{equation}
Do phương sai thực tế khác nhau, $g$ được dùng như chỉ số chuẩn hóa mô tả; chênh lệch
và khoảng tin cậy trên thang gốc vẫn là kết quả ưu tiên để diễn giải.

\subsection{Xây dựng giả thuyết thống kê}
\label{subsec:two-sample-hypotheses}

Với từng RQ, giả thuyết hai phía là
\begin{equation}
  H_0:\mu_N=\mu_A,
  \qquad
  H_1:\mu_N\ne\mu_A.
  \label{eq:two-sample-hypotheses}
\end{equation}
Kiểm định hai phía được chọn vì nghiên cứu không áp đặt trước chiều khác biệt. Ở mức
$\alpha=0{,}05$, $H_0$ bị bác bỏ khi p-value đã hiệu chỉnh nhỏ hơn $\alpha$. Không
bác bỏ $H_0$ chỉ có nghĩa dữ liệu chưa cung cấp đủ bằng chứng cho khác biệt, không phải
chứng minh hai trung bình bằng nhau. Vì có ba kiểm định chính, phương pháp Holm được
dùng để kiểm soát sai lầm loại I theo họ giả thuyết; p-value thô và p-value Holm được
báo cáo riêng.

\subsection{Giả định và chẩn đoán}
\label{subsec:two-sample-assumptions}

Thứ nhất, hai mẫu phải độc lập giữa và trong nhóm. Nhãn Normal và Attack không giao
nhau, nhưng NSL-KDD không cung cấp định danh phiên hoặc máy nguồn đủ để kiểm tra các
kết nối có thuộc cùng một hoạt động hay không. Welch không khắc phục được phụ thuộc
tiềm ẩn này, vì vậy suy luận được xem là khám phá có điều kiện trên cấu trúc dữ liệu.

Thứ hai, phân phối trung bình mẫu cần được xấp xỉ hợp lý. Histogram/ECDF mô tả toàn bộ
phân bố; boxplot cho thấy trung vị, IQR và ngoại lệ; Q--Q plot đối chiếu các phân vị
quan sát với phân vị chuẩn. Shapiro--Wilk trên tối đa 5.000 quan sát mỗi nhóm chỉ là
chẩn đoán bổ sung, không phải quy tắc lựa chọn duy nhất. Cỡ mẫu lớn làm trung bình mẫu
ổn định hơn, nhưng không biến các quan sát gốc thành phân phối chuẩn.

Thứ ba, Brown--Forsythe (Levene với tâm là trung vị) được dùng để đánh giá chênh lệch
phương sai. Welch không yêu cầu phương sai bằng nhau, song vẫn nhạy với phụ thuộc và
các phân bố quá dị biệt. Hai biến tỷ lệ bị chặn và có khối lượng tại 0 hoặc 1;
\texttt{count} là biến đếm lệch phải và có ngoại lệ. Do đó, Mann--Whitney U và
Yuen--Welch với trung bình cắt ngọn 20\% được dùng như phân tích độ nhạy. Chúng không
kiểm định đúng cùng tham số với trung bình gốc nên không thay thế kết quả chính.

\subsection{Quy trình thực hiện}
\label{subsec:two-sample-procedure}

Dữ liệu đầu vào là KDDTrain+ đã làm sạch, còn KDDTest+ được giữ riêng cho đánh giá mô
hình. Nhóm kiểm tra miền giá trị, kiểu dữ liệu, thiếu, NaN và vô hạn; sau đó xác nhận
\texttt{label\_binary}=0 tương ứng Normal và mọi nhãn DoS, Probe, R2L, U2R tương ứng
Attack. Không sử dụng SMOTE, StandardScaler hoặc bản ghi tổng hợp và không loại ngoại
lệ để làm thay đổi p-value.

Trên mỗi đặc trưng, quy trình lần lượt: (i) tính thống kê mô tả và CI trung bình từng
nhóm; (ii) khảo sát boxplot, ECDF, Q--Q plot, skewness, khối lượng tại biên và tỷ số
phương sai; (iii) chạy Welch hai phía trên toàn bộ quan sát hợp lệ; (iv) tính CI của
chênh lệch và Hedges $g$; (v) hiệu chỉnh Holm ba p-value chính; (vi) đối chiếu
Mann--Whitney U và Yuen--Welch; (vii) diễn giải theo hướng dữ liệu, không theo quan hệ
nhân quả. Chỉ ECDF và Q--Q plot được lấy mẫu cố định lần lượt tối đa 5.000 và 1.000
quan sát/nhóm với \texttt{random\_state}=18015; kiểm định dùng toàn bộ dữ liệu.

\subsection{Công cụ triển khai}
\label{subsec:two-sample-tools}

Pandas và NumPy được dùng để đọc, kiểm tra và tổng hợp dữ liệu; SciPy thực hiện các
phân phối, Brown--Forsythe, Mann--Whitney và Yuen--Welch; Statsmodels hiệu chỉnh Holm.
Matplotlib và Seaborn tạo hình PNG 300 DPI và SVG. Mã nguồn, kiểm thử và bảng kết quả
được lưu trong thư mục \texttt{two\_sample\_testing}; phần kết quả được trình bày tại
Mục~\ref{sec:two-sample-results}.

